% ============================================================
% SCHEDE PG AGGIUNTIVE — CASA GRIMALDI
% Personaggi per tavolo alternativo / quinta sessione
% ============================================================

\chapter{Schede PG Aggiuntive --- Casa Grimaldi}
\index{schede PG!Grimaldi!aggiuntivi}

% ════════════════════════════════════════════════════════════
% PG 6 — ANDRIOT DE MONACO detto IL FALCO
% @rpg-schema: <https://rpg-schema.org/instances/genova1507/pg-andriot-monaco>
%   a fitd:Character ;
%   rdfs:label "Andriot de Monaco detto il Falco" ;
%   fitd:crew <.../crew-grimaldi> ;
%   fitd:role "Il Corsaro / Il Braccio Armato di Monaco" ;
%   fitd:action [ fitd:actionName "Hunt" ; fitd:dots 3 ] ;
%   fitd:action [ fitd:actionName "Skirmish" ; fitd:dots 3 ] ;
%   fitd:action [ fitd:actionName "Prowl" ; fitd:dots 3 ] .
% ════════════════════════════════════════════════════════════

\pgheader[GrimaldiPurple]{Andriot de Monaco detto il Falco}{Il Corsaro / Il Braccio Armato di Monaco}{Grimaldi}
\pgquote{Una lettera di corsa del principe di Monaco vale quanto una lettera di corsa del re di Francia. Chiedetelo a chi ho affondato.}

\section*{Background \& Motivazione}

\begin{rulebox}{}
  \textbf{Background:} Trentanove anni, monegasco di nascita, corsaro per vocazione.
  Opera con lettera di corsa firmata da Rainier Grimaldi — il che lo rende
  pirata o privateer a seconda di chi lo cattura. Ha affondato tre navi
  genovesi per conto dei francesi (quando Monaco aveva bisogno dei francesi),
  due navi francesi per conto di mercanti genovesi (quando Monaco aveva bisogno
  dei genovesi), e una nave spagnola per ragioni che preferisce non spiegare.
  È a Genova perché Rainier lo ha convocato. È l'unico della famiglia con
  esperienza diretta di operazioni militari marittime.

  \medskip\noindent\textcolor{GrimaldiPurple}{\textbf{Motivazione:}} Monaco libera
  significa rotte commerciali meno tassate. Rotte meno tassate significa
  prede più ricche. La catena logica è lineare.
\end{rulebox}

\section*{Attributi \& Azioni}


\begin{attributoblock}[GrimaldiPurple]{INSIGHT — Mente}
  \azione{Caccia}{Hunt}{3}{Segue navi, merci, persone — su terra e in mare}
  \azione{Studio}{Study}{1}{Carte nautiche, venti, rotte — nient'altro}
  \azione{Osserva}{Survey}{2}{Valuta pericoli con l'istinto di chi ha sbagliato raramente}
  \azione{Ingegno}{Tinker}{2}{Armi, attrezzi navali, trappole improvvisate}
\end{attributoblock}


\begin{attributoblock}[GrimaldiPurple]{PROWESS — Corpo}
  \azione{Finezza}{Finesse}{2}{Velocità, agilità — si muove come su un ponte in tempesta}
  \azione{Ombra}{Prowl}{3}{Bordate notturne, infiltrazioni costiere, sparizioni rapide}
  \azione{Combattimento}{Skirmish}{3}{Spada corta e pistola da bordo — combatte sporco}
  \azione{Forza bruta}{Wreck}{2}{Sfonda portelli, abbatte alberi, rompe catene}
\end{attributoblock}


\begin{attributoblock}[GrimaldiPurple]{RESOLVE — Spirito}
      \azione{Percezione}{Attune}{1}{Istinto per l'imboscata — navale e terrestre}
    \azione{Comando}{Command}{3}{I marinai lo seguono perché li ha portati a casa ogni volta}
    \azione{Relazioni}{Consort}{2}{Marinai, portuali, taverne di porto, informatori costieri}
    \azione{Persuasione}{Sway}{1}{Diretto. Funziona con chi non ha alternative.}
  \end{attributoblock}

\medskip
\begin{pgclockbox}{Stress \hfill \clock{9}{0}}
  \small\textit{Traumi: $\square$ Freddo \quad $\square$ Duro \quad $\square$ Ossessionato \quad $\square$ Paranoico \quad $\square$ Irrazionale \quad $\square$ Spezzato}
\end{pgclockbox}

\section*{Abilità Speciale}

% @rpg-schema: <.../pg-andriot-monaco>
%   fitd:specialAbility [ rdfs:label "Lettera di Corsa" ] .
\begin{grimaldibox}{Lettera di Corsa}
  Andriot porta una lettera di corsa firmata da Rainier Grimaldi che lo autorizza
  ad agire in nome del Principato di Monaco. Può presentarla per rivendicare
  autorità su qualsiasi azione in mare o sulla costa genovese — incluso il diritto
  di abbordaggio e sequestro di navi francesi in porto.

  \medskip\noindent\textit{Sempre disponibile. Quando presentata, i PNG marittimi
  e i capitani di porto non possono opporsi senza violare il diritto internazionale
  consuetudinario. I francesi possono ignorarla — ma devono motivare il perché.}
\end{grimaldibox}

\section*{Legami}

\begin{table}[h]
\centering\renewcommand{\arraystretch}{1.4}
\begin{tabularx}{\linewidth}{@{}L{2.6cm} X@{}}
  \toprule\rowcolor{GrimaldiPurple}
  \textcolor{white}{\textbf{PG}} & \textcolor{white}{\textbf{Legame}} \\
  \midrule
  Rainier & Il mio principe. Non sempre d'accordo con lui. Sempre fedele a Monaco. \\
  \rowcolor{GrimaldiPurple!8}
  Ser Battista & Il capitano è il miglior soldato che abbia mai visto in terra. In mare avrei da ridire. Ma è una questione teorica. \\
  Costantino & Il mercante sa dove sono le navi e cosa trasportano. Per me è più utile di qualsiasi alleato armato. \\
  \rowcolor{GrimaldiPurple!8}
  Torquato & L'oste. Mi ha salvato la vita una volta in un vicolo di Soziglia. Non gli ho ancora chiesto perché. \\
  \bottomrule
\end{tabularx}
\end{table}

\section*{Equipaggiamento}
\begin{goldlist}
  \item Lettera di corsa originale firmata da Rainier Grimaldi — con sigillo
  \item Spada da bordo e pistola a pietra focaia (un colpo solo)
  \item Mappa del porto di Genova con i fondali reali e le rotte di fuga costiere
  \item Chiave di un magazzino a Sampierdarena con sei marinai armati ad aspettare
\end{goldlist}


\begin{secretbox}
  \textbf{Segreto:} Andriot ha una seconda lettera — non di Rainier, ma di un
  armatore francese di Marsiglia che gli ha promesso immunità e un carico
  di merci se consegna informazioni sui movimenti della flotta genovese
  durante la rivolta. Non ha ancora deciso se accettare. Il carico vale
  abbastanza da comprare una nave propria. Una nave propria significherebbe
  smettere di dipendere dalla lettera di corsa di Rainier.

  \medskip\noindent\textcolor{SecretRed}{\textbf{Agenda:}}
  \textit{Completare la notte con Monaco riconosciuta \textbf{e} con abbastanza
  informazioni vendibili da usare come uscita di emergenza indipendente.
  Non tradire Rainier — ma costruirsi una via che non passi per Rainier.}
\end{secretbox}

\clearpage
% ════════════════════════════════════════════════════════════
% PG 7 — TORQUATO GRIMALDI detto L'OSTE
% @rpg-schema: <https://rpg-schema.org/instances/genova1507/pg-torquato-grimaldi>
%   a fitd:Character ;
%   rdfs:label "Torquato Grimaldi detto l'Oste" ;
%   fitd:crew <.../crew-grimaldi> ;
%   fitd:role "L'Oste dei Caruggi / La Radice Genovese" ;
%   fitd:action [ fitd:actionName "Consort" ; fitd:dots 4 ] ;
%   fitd:action [ fitd:actionName "Survey" ; fitd:dots 3 ] ;
%   fitd:action [ fitd:actionName "Hunt" ; fitd:dots 2 ] .
% ════════════════════════════════════════════════════════════

\pgheader[GrimaldiPurple]{Torquato Grimaldi detto l'Oste}{L'Oste dei Caruggi / La Radice Genovese}{Grimaldi}
\pgquote{Monaco è il Principato. Genova è casa mia. Tengo entrambe le cose in testa contemporaneamente. Non è difficile quanto sembra.}

\section*{Background \& Motivazione}

\begin{rulebox}{}
  \textbf{Background:} Cinquantasei anni, nato a Genova da padre monegasco e madre
  genovese. Gestisce una locanda nel sestiere di Soziglia da trent'anni —
  ufficialmente di proprietà di un prestanome genovese, di fatto un asset
  dei Grimaldi per raccogliere informazioni, ospitare agenti in transito,
  e mantenere una presenza stabile nella città. Torquato non è mai andato
  a Monaco. È genovese in tutto tranne il cognome e il padrone.
  Conosce Genova come le proprie tasche — i caruggi, le voci, i rancori
  e le fedeltà di quartiere che nessun nobile capisce davvero.

  \medskip\noindent\textcolor{GrimaldiPurple}{\textbf{Motivazione:}} È stanco
  dei francesi. Non per ragioni politiche astratte — perché i loro soldati
  bevono senza pagare e rompono i tavoli. E perché sua figlia ha smesso
  di uscire da quando ci sono le pattuglie notturne.
\end{rulebox}

\section*{Attributi \& Azioni}


\begin{attributoblock}[GrimaldiPurple]{INSIGHT — Mente}
  \azione{Caccia}{Hunt}{2}{Trova chiunque a Genova — sa chi frequenta cosa}
  \azione{Studio}{Study}{1}{Sa leggere, far di conto, nient'altro di accademico}
  \azione{Osserva}{Survey}{3}{Trent'anni di locanda insegnano a leggere le persone}
  \azione{Ingegno}{Tinker}{1}{Cucina, riparazioni domestiche, trappole semplici}
\end{attributoblock}


\begin{attributoblock}[GrimaldiPurple]{PROWESS — Corpo}
  \azione{Finezza}{Finesse}{1}{Mani da cuoco e da lottatore — usa entrambe}
  \azione{Ombra}{Prowl}{2}{Si muove nei caruggi come a casa sua — perché è casa sua}
  \azione{Combattimento}{Skirmish}{2}{Mazza, coltello da cucina, qualsiasi cosa sottomano}
  \azione{Forza bruta}{Wreck}{2}{Robusto. Ha buttato fuori ubriachi per trent'anni.}
\end{attributoblock}


\begin{attributoblock}[GrimaldiPurple]{RESOLVE — Spirito}
      \azione{Percezione}{Attune}{2}{Sente quando una conversazione porta guai}
    \azione{Comando}{Command}{1}{Autorità di quartiere — funziona con i vicini}
    \azione{Relazioni}{Consort}{4}{Conosce tutti in Soziglia e Portoria — e i loro debiti}
    \azione{Persuasione}{Sway}{2}{Convince con il bicchiere e il buon senso}
  \end{attributoblock}

\medskip
\begin{pgclockbox}{Stress \hfill \clock{9}{0}}
  \small\textit{Traumi: $\square$ Freddo \quad $\square$ Duro \quad $\square$ Ossessionato \quad $\square$ Paranoico \quad $\square$ Irrazionale \quad $\square$ Spezzato}
\end{pgclockbox}

\section*{Abilità Speciale}

% @rpg-schema: <.../pg-torquato-grimaldi>
%   fitd:specialAbility [ rdfs:label "La Locanda Sa Tutto" ] .
\begin{grimaldibox}{La Locanda Sa Tutto}
  La locanda di Torquato è frequentata da persone di ogni ceto e fazione.
  Qualsiasi cosa accada a Genova, qualcuno che lo sa passa dalla sua porta.
  Una volta per sessione, Torquato può dichiarare di aver sentito qualcosa
  di rilevante nella locanda nelle ultime ore: un piano, una posizione,
  un segreto circolato tra i clienti.

  \medskip\noindent\textit{Nessun tiro. L'informazione è accurata ma frammentaria
  — può essere incompleta o decontestualizzata.}
\end{grimaldibox}

\section*{Legami}

\begin{table}[h]
\centering\renewcommand{\arraystretch}{1.4}
\begin{tabularx}{\linewidth}{@{}L{2.6cm} X@{}}
  \toprule\rowcolor{GrimaldiPurple}
  \textcolor{white}{\textbf{PG}} & \textcolor{white}{\textbf{Legame}} \\
  \midrule
  Rainier & Il mio signore. Lo rispetto. Non lo conosco — non ci siamo mai parlati direttamente. \\
  \rowcolor{GrimaldiPurple!8}
  Costantino & Il mercante greco porta i clienti migliori alla mia locanda. Facciamo buoni affari. Non ci fidiamo l'uno dell'altro. \\
  Serafina & La figlia del signore è venuta alla locanda una volta, da sola, travestita. Non ho detto niente a Rainier. \\
  \rowcolor{GrimaldiPurple!8}
  Andriot & Il corsaro. Gli ho salvato la vita in un vicolo. Mi deve un favore. È il favore più grande che abbia mai tenuto in serbo. \\
  \bottomrule
\end{tabularx}
\end{table}

\section*{Equipaggiamento}
\begin{goldlist}
  \item Chiavi di undici locali nel sestiere di Soziglia — cantine, magazzini, uscite secondarie
  \item Lista di chi deve cosa alla locanda — debiti di trenta anni, alcuni inestimabili
  \item Mappa mentale di tutti i tunnel di drenaggio tra Soziglia e il porto (non scritta)
  \item Botte di vino adulterato — effetto soporifero in dosi elevate, «per le emergenze»
\end{goldlist}


\begin{secretbox}
  \textbf{Segreto:} Torquato sa dove è nascosta la figlia di Rainier — Serafina
  è venuta alla locanda una settimana fa chiedendo di affittare una stanza
  con un nome falso. Torquato l'ha riconosciuta, ha preso i soldi, e non ha
  detto nulla. Ora sa che Serafina si sta preparando a fare qualcosa senza
  il permesso del padre. Non sa cosa. Sa che se glielo chiede, potrebbe
  dover scegliere tra Serafina e Rainier.

  \medskip\noindent\textcolor{SecretRed}{\textbf{Agenda:}}
  \textit{Tenere la figlia del signore al sicuro senza che Rainier lo sappia,
  senza che Serafina sappia che lui sa, e senza dover scegliere tra i due.
  Torquato è convinto di poter gestire tutte e tre le cose. Ha sessant'anni
  di esperienza nel gestire situazioni impossibili. Questa volta potrebbe
  avere torto.}
\end{secretbox}
